% GENERATED FILE. Edit canonical lessons, then run npm run generate:books.
% canonical-source-hash: fnv1a64:3fe406bcd1b89269
% canonical-lessons: FA-C311-manba, FA-C311-be-gofte, FA-C311-be-naql-az, FA-C311-joziyat, FA-C311-khabargozari

\chapter{Sources and details}
\label{ch:fa-sources-and-details}
\hlchaptermodality{311}

\begin{chapteropening}
\textbf{By the end of this chapter:} \emph{I can say where news came from: a source, according to someone, quoting someone, the details and a news agency.}

Complete the last lesson of chapter 311: I can say where news came from: a source, according to someone, quoting someone, the details and a news agency.
\end{chapteropening}

\section[\texorpdfstring{\fa{منبع} (manba')}{منبع (manba')}]{\fa{منبع} (manba') — a source}

\label{lesson:FA-C311-manba}



\begin{quote}
Before the new one: say the Persian for to let someone know, then the Persian for to inform.
\end{quote}

\subsection*{You'll want to know: \fa{منبع}}
\textbf{\fa{منبع}} (\emph{manba'}) — \textquotedblleft{}a source\textquotedblright{}.

\begin{quote}
\fa{این} \fa{خبر} \fa{از} \fa{یک} \fa{منبع} \fa{مطمئن} \fa{بود}.
\end{quote}

\begin{quote}
\emph{In khabar az yek manba'-e motma'en bud.}
\end{quote}

\begin{quote}
This news came from a reliable source.
\end{quote}

Broken plural \emph{manâbe'}. Its first meaning is a spring of water.

\begin{cousinweb}
\emph{manba'} is borrowed from Arabic \emph{manba'}, \textquotedblleft{}spring, source\textquotedblright{} (EtymDB).
\end{cousinweb}

\begin{grammarlens}[title={What you've built}]
One more word for telling, explaining and describing at B1.
\end{grammarlens}

\subsection*{Your turn}
\begin{itemize}
\raggedright
  \item \emph{Say it:} \emph{manba'}
  \item \emph{Say it:} \emph{manba'}, once more
  \item \emph{Say it:} a sentence of your own with \emph{manba'}
  \item \emph{Recall:} say the Persian for to let someone know, then the Persian for to inform, then say \emph{manba'} again
\end{itemize}

\begin{tcolorbox}[breakable,colback=teal!4,colframe=teal!35!black,title={Before you move on}]
What does \fa{منبع} mean? (\textquotedblleft{}A source\textquotedblright{}.) Say it once more.
\end{tcolorbox}
\section[\texorpdfstring{\fa{به} \fa{گفته} (be gofte-ye)}{به گفته (be gofte-ye)}]{\fa{به} \fa{گفته} (be gofte-ye) — according to}

\label{lesson:FA-C311-be-gofte}



\begin{quote}
Before the new one: say the Persian for to inform, then the Persian for a source.
\end{quote}

\subsection*{You'll want to know: \fa{به} \fa{گفته}}
\textbf{\fa{به} \fa{گفته}} (\emph{be gofte-ye}) — \textquotedblleft{}according to\textquotedblright{}.

\begin{quote}
\fa{به} \fa{گفته} \fa{پلیس،} \fa{کسی} \fa{آسیب} \fa{ندید}.
\end{quote}

\begin{quote}
\emph{Be gofte-ye polis, kasi âsib nadid.}
\end{quote}

\begin{quote}
According to the police, nobody was hurt.
\end{quote}

It takes the ezafe: \emph{be gofte-ye u} \textquotedblleft{}according to him\textquotedblright{}. \emph{gofte} \textquotedblleft{}what was said\textquotedblright{} is the participle of \emph{goftan}.

\begin{cousinweb}
\emph{be gofte-ye} is built from \emph{be} \textquotedblleft{}to\textquotedblright{} and \emph{gofte} \textquotedblleft{}what was said\textquotedblright{}, the participle of \emph{goftan} \textquotedblleft{}to say\textquotedblright{}. All of it is native Persian.
\end{cousinweb}

\begin{grammarlens}[title={What you've built}]
One more word for telling, explaining and describing at B1.
\end{grammarlens}

\subsection*{Your turn}
\begin{itemize}
\raggedright
  \item \emph{Say it:} \emph{be gofte-ye}
  \item \emph{Say it:} \emph{be gofte-ye}, once more
  \item \emph{Say it:} a sentence of your own with \emph{be gofte-ye}
  \item \emph{Recall:} say the Persian for to inform, then the Persian for a source, then say \emph{be gofte-ye} again
\end{itemize}

\begin{tcolorbox}[breakable,colback=teal!4,colframe=teal!35!black,title={Before you move on}]
What does \fa{به} \fa{گفته} mean? (\textquotedblleft{}According to\textquotedblright{}.) Say it once more.
\end{tcolorbox}
\section[\texorpdfstring{\fa{به} \fa{نقل} \fa{از} (be naql az)}{به نقل از (be naql az)}]{\fa{به} \fa{نقل} \fa{از} (be naql az) — quoting, citing}

\label{lesson:FA-C311-be-naql-az}



\begin{quote}
Before the new one: say the Persian for a source, then the Persian for according to.
\end{quote}

\subsection*{You'll want to know: \fa{به} \fa{نقل} \fa{از}}
\textbf{\fa{به} \fa{نقل} \fa{از}} (\emph{be naql az}) — \textquotedblleft{}quoting, citing\textquotedblright{}.

\begin{quote}
\fa{روزنامه} \fa{به} \fa{نقل} \fa{از} \fa{وزیر} \fa{نوشت} \fa{که} \fa{قیمت} \fa{ها} \fa{بالا} \fa{میرود}.
\end{quote}

\begin{quote}
\emph{Ruznâme be naql az vazir nevesht ke qeymat-hâ bâlâ miravad.}
\end{quote}

\begin{quote}
Quoting the minister, the newspaper wrote that prices were going up.
\end{quote}

\emph{naql} \textquotedblleft{}transfer, relaying\textquotedblright{} is the same noun as in \emph{naql kardan} \textquotedblleft{}to relate\textquotedblright{}.

\begin{cousinweb}
\emph{naql} is borrowed from Arabic \emph{naql}, \textquotedblleft{}transport, transfer\textquotedblright{} (EtymDB); \emph{be} \textquotedblleft{}to\textquotedblright{} and \emph{az} \textquotedblleft{}from\textquotedblright{} are inherited Persian.
\end{cousinweb}

\begin{grammarlens}[title={What you've built}]
One more word for telling, explaining and describing at B1.
\end{grammarlens}

\subsection*{Your turn}
\begin{itemize}
\raggedright
  \item \emph{Say it:} \emph{be naql az}
  \item \emph{Say it:} \emph{be naql az}, once more
  \item \emph{Say it:} a sentence of your own with \emph{be naql az}
  \item \emph{Recall:} say the Persian for a source, then the Persian for according to, then say \emph{be naql az} again
\end{itemize}

\begin{tcolorbox}[breakable,colback=teal!4,colframe=teal!35!black,title={Before you move on}]
What does \fa{به} \fa{نقل} \fa{از} mean? (\textquotedblleft{}Quoting, citing\textquotedblright{}.) Say it once more.
\end{tcolorbox}
\section[\texorpdfstring{\fa{جزئیات} (joz'iyât)}{جزئیات (joz'iyât)}]{\fa{جزئیات} (joz'iyât) — details, the details}

\label{lesson:FA-C311-joziyat}



\begin{quote}
Before the new one: say the Persian for according to, then the Persian for quoting.
\end{quote}

\subsection*{You'll want to know: \fa{جزئیات}}
\textbf{\fa{جزئیات}} (\emph{joz'iyât}) — \textquotedblleft{}details, the details\textquotedblright{}.

\begin{quote}
\fa{پلیس} \fa{جزئیات} \fa{بیشتری} \fa{نگفت}.
\end{quote}

\begin{quote}
\emph{Polis joz'iyât-e bishtari nagoft.}
\end{quote}

\begin{quote}
The police gave no further details.
\end{quote}

An Arabic plural built on \emph{joz'i} \textquotedblleft{}partial, minor\textquotedblright{} (from \emph{joz'} \textquotedblleft{}a part\textquotedblright{}), written with the hamza seat \fa{ئ}.

\begin{cousinweb}
\emph{joz'iyât} is borrowed from Arabic \emph{juz'iyyāt}, \textquotedblleft{}details\textquotedblright{} (EtymDB), a plural built on \emph{juz'} \textquotedblleft{}a part\textquotedblright{}.
\end{cousinweb}

\begin{grammarlens}[title={What you've built}]
One more word for telling, explaining and describing at B1.
\end{grammarlens}

\subsection*{Your turn}
\begin{itemize}
\raggedright
  \item \emph{Say it:} \emph{joz'iyât}
  \item \emph{Say it:} \emph{joz'iyât}, once more
  \item \emph{Say it:} a sentence of your own with \emph{joz'iyât}
  \item \emph{Recall:} say the Persian for according to, then the Persian for quoting, then say \emph{joz'iyât} again
\end{itemize}

\begin{tcolorbox}[breakable,colback=teal!4,colframe=teal!35!black,title={Before you move on}]
What does \fa{جزئیات} mean? (\textquotedblleft{}Details, the details\textquotedblright{}.) Say it once more.
\end{tcolorbox}
\section[\texorpdfstring{\fa{خبرگزاری} (khabargozâri)}{خبرگزاری (khabargozâri)}]{\fa{خبرگزاری} (khabargozâri) — a news agency}

\label{lesson:FA-C311-khabargozari}



\begin{quote}
Before the new one: say the Persian for quoting, then the Persian for details.
\end{quote}

\subsection*{You'll want to know: \fa{خبرگزاری}}
\textbf{\fa{خبرگزاری}} (\emph{khabargozâri}) — \textquotedblleft{}a news agency\textquotedblright{}.

\begin{quote}
\fa{این} \fa{خبر} \fa{را} \fa{اول} \fa{یک} \fa{خبرگزاری} \fa{خارجی} \fa{داد}.
\end{quote}

\begin{quote}
\emph{In khabar râ avval yek khabargozâri-ye khâreji dâd.}
\end{quote}

\begin{quote}
A foreign news agency gave this news first.
\end{quote}

\emph{khabar} \textquotedblleft{}news\textquotedblright{} + \emph{gozâri} \textquotedblleft{}rendering\textquotedblright{}, from the stem \emph{gozâr-}.

\begin{cousinweb}
\emph{khabargozâri} joins \emph{khabar} \textquotedblleft{}news\textquotedblright{}, borrowed from Arabic \emph{khabar} (EtymDB), and the native Persian \emph{gozâri} \textquotedblleft{}rendering\textquotedblright{}, from the stem \emph{gozâr-}.
\end{cousinweb}

\begin{grammarlens}[title={What you've built}]
One more word for telling, explaining and describing at B1.
\end{grammarlens}

\subsection*{Your turn}
\begin{itemize}
\raggedright
  \item \emph{Say it:} \emph{khabargozâri}
  \item \emph{Say it:} \emph{khabargozâri}, once more
  \item \emph{Say it:} a sentence of your own with \emph{khabargozâri}
  \item \emph{Recall:} say the Persian for quoting, then the Persian for details, then say \emph{khabargozâri} again
\end{itemize}

\begin{tcolorbox}[breakable,colback=teal!4,colframe=teal!35!black,title={Before you move on}]
What does \fa{خبرگزاری} mean? (\textquotedblleft{}A news agency\textquotedblright{}.) Say it once more.
\end{tcolorbox}
